\lecture{1}{Fri 12 Jan 2024 13:15}{third lecture}

This is one of the most important theorems in this course. Also very important in probability.

\begin{theorem}[Borel Cantelli Lemma]
	\label{thm:BorelCantelliLem}
	
	Suppose $(X, \mathcal{A}, \mu)$ is a measure space, let $E_k$, $k \ge 1$, be a sequence of sets in $\mathcal{A}$ with $\sum_{k = 1}^\infty  \mu(E_k) < \infty $. Then $\mu(\limsup E_k) = 0$.
\end{theorem}
\begin{proof}
	Let $\limsup E_k = \bigcap_{m = 1} ^\infty  \bigcup_{n = m} ^\infty E_n = \bigcap_{m = 1}^\infty {B_m}  $.
	Note $B_m$ are decreasing, also
	\[
		\mu(B_1) = \mu\left( \bigcup_{n = 1} ^\infty  E_n \right) \le \sum_{n = 1}^\infty  \mu(E_n) < \infty 
	.\] 
	Apply continuity from above,
	\[
		\mu\left( \limsup E_k \right)  = \mu \left( \lim B_k \right) = \lim \mu\left( B_k \right) 
	.\] 
	But
	\[
		\mu(B_k) = \mu\left( \bigcup_{n = k} ^\infty  E_n \right) \le \sum_{n = k}^\infty  \mu(E_n) \to 0
	.\] 
\end{proof}

